% Source: https://artofproblemsolving.com/community/c3h1322461_bees_problem
% Stable identity: aops-topic-1322461-q01
Let there be $n$ bees. The fractions going to the trees are $\dfrac{1}{5}$, $\dfrac{1}{3}$ and $3\left(\dfrac{1}{3}-\dfrac{1}{5}\right)=\dfrac{2}{5}$, a total of $\dfrac{3+5+6}{15}=\dfrac{14}{15}$. So the one remaining bee is $\dfrac{1}{15}$ of the swarm, and $n=\boxed{15}$. (The poster's $\dfrac{9}{15}$ omitted the factor $3$ in the acorn-tree group.)
